\documentclass[UTF8,11pt,a4paper]{ctexart}
\usepackage[a4paper,left=24mm,right=24mm,top=24mm,bottom=23mm,headheight=15pt]{geometry}
\usepackage{xcolor,graphicx,booktabs,longtable,tabularx,array,multirow}
\usepackage{tikz}
\usetikzlibrary{arrows.meta,positioning,fit,shapes.geometric,calc,backgrounds}
\usepackage{enumitem,listings,fancyhdr,hyperref,titlesec,lastpage}
\definecolor{navy}{HTML}{173B57}
\definecolor{blue}{HTML}{247BA0}
\definecolor{cyan}{HTML}{5BC0BE}
\definecolor{ink}{HTML}{23313D}
\definecolor{muted}{HTML}{657786}
\definecolor{paper}{HTML}{F4F7F9}
\definecolor{orange}{HTML}{F2994A}
\hypersetup{colorlinks=true,linkcolor=navy,urlcolor=blue,pdfauthor={Slurm AI Agent 项目组},pdftitle={Slurm AI Agent 设计文档}}
\setlength{\parindent}{2em}
\setlength{\parskip}{0.25em}
\setlength{\emergencystretch}{3em}
\setlist{nosep,leftmargin=2em}
\setlist[itemize]{label=\textcolor{blue}{\small\textbullet}}
\pagestyle{fancy}
\fancyhf{}
\fancyhead[L]{\small\textcolor{muted}{Slurm AI Agent 设计文档}}
\fancyhead[R]{\small\textcolor{muted}{USTC 107 算力平台}}
\fancyfoot[C]{\small\textcolor{muted}{第 \thepage\ 页 / 共 \pageref*{LastPage} 页}}
\renewcommand{\headrulewidth}{0.4pt}
\titleformat{\section}{\Large\bfseries\color{navy}}{\thesection}{0.7em}{}[\vspace{-0.2em}\color{cyan}\titlerule]
\titleformat{\subsection}{\large\bfseries\color{blue}}{\thesubsection}{0.6em}{}
\titleformat{\subsubsection}{\normalsize\bfseries\color{ink}}{\thesubsubsection}{0.5em}{}
\lstset{basicstyle=\ttfamily\small,backgroundcolor=\color{paper},frame=single,rulecolor=\color{cyan},breaklines=true,columns=fullflexible,keepspaces=true,showstringspaces=false,aboveskip=0.7em,belowskip=0.7em}
\newcommand{\code}[1]{\texttt{\small #1}}
\newcommand{\tagbox}[1]{\colorbox{paper}{\textcolor{navy}{\strut\quad #1\quad}}}

\begin{document}
\begin{titlepage}
  \pagecolor{navy}
  \color{white}
  \vspace*{22mm}
  {\Huge\bfseries Slurm AI Agent\par}
  \vspace{5mm}
  {\LARGE 面向 USTC 107 算力平台的智能作业助手\par}
  \vspace{12mm}
  {\color{cyan}\rule{0.76\textwidth}{1.5pt}\par}
  \vspace{12mm}
  {\fontsize{30}{36}\selectfont\bfseries 设计文档\par}
  \vspace{10mm}
  {\Large 从自然语言意图到可验证 Slurm 作业的闭环设计\par}
  \vfill
  \begin{tabular}{@{}ll@{}}
    文档版本：& V1.0\\[2mm]
    项目版本：& 基于 2026-09-03 源码快照\\[2mm]
    技术栈：& FastAPI / OpenAI SDK / Slurm REST API / WebSocket\\[2mm]
    部署目标：& USTC 107 算力平台登录节点
  \end{tabular}
  \vspace{18mm}
\end{titlepage}
\nopagecolor
\color{ink}

\begin{abstract}
Slurm AI Agent 是一套部署在共享 HPC 登录节点上的智能作业工作台。系统将项目文件、隔离环境、资源选择、作业提交、运行监控、日志分析和平台知识问答整合到一个 Web 界面，并通过大模型 Function Calling 将自然语言请求转换为可审计、可校验的工具调用。其设计重点不在于让模型直接执行任意命令，而在于把大模型限定在“理解与编排”层，把权限、资源、路径和提交过程固定在确定性后端中，形成从项目准备到结果复盘的安全闭环。
\end{abstract}

\tableofcontents
\clearpage

\section{项目概述}

\subsection{背景与痛点}
Slurm 是科研计算和深度学习训练中常用的集群调度系统，但对初学者并不友好：分区、账户、QoS、TRES、作业脚本、环境管理和日志路径，任何一处都可能成为第一次提交的障碍。传统命令行流程还把任务上下文分散在文件系统、终端、调度队列和文档之间。归纳起来，痛点集中在五个方面：

\begin{enumerate}
  \item \textbf{资源配置易错：}参数组合繁多，写错即提交失败，失败原因还常常不直观。
  \item \textbf{依赖环境易坏：}版本冲突、来源混装、误装进 base 环境，排查成本高。
  \item \textbf{上下文分散：}代码、数据、环境、命令、日志和文档各在一处，任务被打断后难以接续。
  \item \textbf{失败难以复盘：}排队原因无解释，报错堆栈不友好，日志位置靠记忆。
  \item \textbf{共享节点安全顾虑：}多用户环境下，开放端口、明文密钥和不受限的命令执行都是风险。
\end{enumerate}

本项目面向 USTC 107 算力平台，把这五个痛点转化为一条项目级工作流逐项应对。用户仍然决定安装哪些依赖、申请多少资源以及何时提交；系统负责采集真实上下文、提供建议、执行校验并持续跟踪结果。这一“人做决策、模型做理解、程序做约束”的职责划分，是整体设计的核心。

\subsection{设计目标：从痛点到对策}
\noindent\begin{tabularx}{\textwidth}{>{\color{muted}\small}p{26mm}>{\bfseries\color{navy}}p{24mm}X}
\toprule
\textbf{痛点} & \textbf{设计目标} & \textbf{设计对策}\\
\midrule
资源配置易错 & 真实可验证 & 资源与作业状态一律来自实时接口；提交前做授权、配额与路径校验，提交后回读 TRES/GRES 核验 GPU 请求是否生效。\\
依赖环境易坏 & 可控安装、可自愈 & 依赖清单由确定性扫描与真实软件源预检生成，用户确认后才执行；安装命令规范化并互斥运行，中断留下的残缺环境可自愈重建。\\
上下文分散 & 上下文一致 & 以项目为根组织代码、数据、环境、会话与作业记录；文件一旦变更即标记依赖分析失效，避免复用过期结论。\\
失败难以复盘 & 结果可追溯 & 提交后的作业注册为观察记录，终态自动聚合 Slurm 状态与日志生成报告；失败中的缺失依赖写回需求，反哺下一轮检查。\\
共享节点安全顾虑 & 安全可控 & 密钥不写入代码、网络默认隧道访问、CLI 白名单；安装、提交等有副作用的操作均由用户明确触发。\\
\bottomrule
\end{tabularx}

贯穿上述对策的手段是自然语言交互：模型负责理解与解释，降低 Slurm 概念的学习成本，但不改变表中任何一条约束。另有一项横切目标是\textbf{可降级}：Embedding 或模型分析失败时，分别退回关键词检索或原始日志报告，保证核心功能在外部服务异常时不中断。

\subsection{设计范围与非目标}
系统覆盖 Web 工作区、智能体、多源 Slurm 能力、依赖规划、作业模板、受控提交、监控分析、文件管理和登录节点终端。不将大模型作为授权主体，不允许其自由拼装写操作 CLI；不替代 Slurm 调度器，也不承诺改变排队优先级；交互式 \code{srun --pty} 等需真实 TTY 的工作流仍由用户在终端完成。

\section{总体设计思路}

\subsection{从“聊天机器人”到“受控执行代理”}
系统把任务拆为三类能力：大模型负责语义理解、步骤规划和结果解释；工具层负责将结构化意图映射到有限能力；确定性后端负责权限、参数、路径、环境和副作用控制。模型的输出不是最终提交脚本，而是结构化资源字段和命令正文。服务端随后生成锁定的脚本前导、工作目录、Conda 激活和日志配置，从而避免提示注入或模型幻觉直接改变调度边界。

\begin{center}
\begin{tikzpicture}[node distance=8mm and 8mm,>=Latex,font=\small,
  box/.style={rounded corners=2mm,draw=blue,very thick,fill=paper,minimum height=10mm,align=center,text width=29mm},
  gate/.style={rounded corners=2mm,draw=orange,very thick,fill=orange!10,minimum height=10mm,align=center,text width=31mm},
  arr/.style={->,thick,color=muted}]
  \node[box] (u) {用户意图\\项目与资源选择};
  \node[box,right=of u] (m) {大模型编排\\Function Calling};
  \node[gate,right=of m] (g) {确定性校验网关\\授权/配额/路径/语法};
  \node[box,right=of g] (s) {Slurm 与项目环境\\执行及状态回读};
  \draw[arr] (u)--(m); \draw[arr] (m)--(g); \draw[arr] (g)--(s);
  \draw[arr,bend left=25] (s) to (m);
  \draw[arr,bend right=27] (m) to (u);
\end{tikzpicture}
\end{center}

\subsection{项目化状态模型}
默认工作空间为 \code{\textasciitilde/projects/<项目名>}。每个项目共享一个 \code{.slurm-agent/conda-env}，可包含多个运行子目录；每个运行目录维护独立对话上下文，同时共享项目依赖。这种模型兼顾实验隔离和依赖复用，也便于把提交脚本、标准输出、错误日志和报告统一放入 \code{logs/}。

\begin{lstlisting}
~/projects/<project>/
|-- activate.sh
|-- <source, data, config>
|-- <run-subdir>/
|-- logs/
|-- runs/
`-- .slurm-agent/
    |-- conda-env/
    |-- dependency-analysis.json
    `-- install.lock
\end{lstlisting}

目录结构背后的共享关系如下图：多个运行子目录各自维护独立对话上下文，但激活同一个项目级 Conda 环境，安装与修改受互斥锁保护。

\begin{center}
\begin{tikzpicture}[node distance=6mm and 10mm, font=\small,
  box/.style={rounded corners=1.5mm,draw=blue,thick,fill=blue!8,minimum height=10mm,align=center,text width=34mm},
  dep/.style={rounded corners=1.5mm,draw=orange,thick,fill=orange!10,minimum height=10mm,align=center,text width=46mm},
  a/.style={->,thick,color=muted}]
  \node[box] (run1) {运行子目录 A\\独立对话会话};
  \node[box,right=of run1] (run2) {运行子目录 B\\独立对话会话};
  \node[dep] (env) at ($(run1.south)!0.5!(run2.south)+(0,-1.1)$) {\code{.slurm-agent/conda-env}\\项目共享依赖（互斥修改）};
  \draw[a] (run1) -- (env);
  \draw[a] (run2) -- (env);
  \begin{scope}[on background layer]
    \node[rounded corners=2mm,draw=navy,thick,fill=paper,fit=(run1)(run2)(env),inner sep=4mm,
      label={[anchor=north west,text=muted,font=\scriptsize]north west:项目目录 \code{\textasciitilde/projects/<项目名>}}] {};
  \end{scope}
\end{tikzpicture}
\end{center}

\section{技术架构}

\subsection{分层架构}
\begin{center}
\begin{tikzpicture}[font=\small,node distance=4mm,
 layer/.style={rounded corners=1.5mm,draw=navy,thick,minimum width=145mm,minimum height=12mm,align=center},
 arr/.style={-{Latex[length=2mm]},thick,color=muted}]
 \node[layer,fill=navy!8] (ui) {表现层：原生 HTML/CSS/JavaScript、SSE 对话、xterm.js 终端、文件与作业面板};
 \node[layer,fill=blue!8,below=of ui] (api) {服务层：FastAPI 路由、项目/会话状态、依赖流程、监控与报告后台线程};
 \node[layer,fill=cyan!12,below=of api] (agent) {智能层：LLMProvider、AgentLoop、19 项 Function Calling 工具、提示词约束};
 \node[layer,fill=orange!10,below=of agent] (core) {领域层：SlurmClient、只读 CLI、知识库、模板引擎、依赖规划、文件传输};
 \node[layer,fill=paper,below=of core] (infra) {基础设施：Slurm REST/CLI、OpenAI 兼容 API、项目文件系统、Conda、Unix Socket/SSH 隧道};
 \draw[arr] (ui)--(api); \draw[arr] (api)--(agent); \draw[arr] (agent)--(core); \draw[arr] (core)--(infra);
\end{tikzpicture}
\end{center}

\subsection{核心组件职责}
\begin{longtable}{>{\raggedright\arraybackslash}p{27mm}>{\raggedright\arraybackslash}p{48mm}>{\raggedright\arraybackslash}p{74mm}}
\toprule
\textbf{组件} & \textbf{实现位置} & \textbf{主要职责}\\
\midrule
\endfirsthead
\toprule
\textbf{组件} & \textbf{实现位置} & \textbf{主要职责}\\
\midrule
\endhead
Web/API 服务 & \path{server/app.py} & 49 个 HTTP/WebSocket 端点，组织项目、模型、仪表盘、依赖、作业、报告、文件等业务流程。\\
智能体循环 & \path{agent/agent_loop.py} & 维护消息历史，接收模型工具调用，逐项执行并回填 tool 消息，直到模型产生最终文本。\\
模型适配 & \path{agent/llm_provider.py} & 使用 OpenAI Python SDK 调用兼容接口，动态读取当前模型，支持工具定义、温度和输出长度控制。\\
工具注册表 & \path{agent/tools_registry.py} & 声明 19 项工具的 JSON Schema，并把工具名路由到 REST、CLI、知识库、模板或项目文件能力。\\
Slurm 客户端 & \path{core/slurm_client.py} & 封装 slurmrestd v0.0.41，请求头认证、401 自动刷新重试、错误规范化、提交与取消。\\
只读 CLI 层 & \path{core/cli_executor.py} & 补足 REST 未覆盖的优先级、fairshare、实时占用、用量报表和账户授权查询。\\
知识库 & \path{core/knowledge_base.py} & Markdown 按标题分块，Embedding 批量向量化、磁盘缓存、余弦检索和关键词降级。\\
依赖规划 & \path{core/dependency_planner.py} & 扫描多种依赖文件与源码 import，版本预检、合并去重、安装计划生成和安装后验证。\\
文件与环境 & \path{core/file_transfer.py} & 工作区创建、安全相对路径、上传限制、项目级 Conda 和跨线程/进程安装锁。\\
终端桥接 & \path{server/terminal.py} & WebSocket 双向字节流连接 PTY/bash，处理尺寸变化、并发限制、Origin 校验与断连清理。\\
\bottomrule
\end{longtable}

\subsection{外部接口与通信协议}
\begin{longtable}{>{\bfseries\color{navy}}p{22mm}>{\raggedright\arraybackslash}p{64mm}>{\raggedright\arraybackslash}p{63mm}}
\toprule
\textbf{接口} & \textbf{协议与地址} & \textbf{说明}\\
\midrule
\endfirsthead
\toprule
\textbf{接口} & \textbf{协议与地址} & \textbf{说明}\\
\midrule
\endhead
模型对话 & OpenAI 兼容 \code{POST /chat/completions} & 基址 \url{https://api.llm.ustc.edu.cn/v1}；携带 tools 时启用 Function Calling；Key 仅取自 \code{LLM\_API\_KEY}。\\
文本向量 & OpenAI 兼容 \code{POST /embeddings} & 同一基址；用于知识库分块向量化与查询向量生成。\\
Slurm REST & \url{http://107.ustc.edu.cn:6820}，前缀 \nolinkurl{/slurm/v0.0.41} & 请求头 \code{X-SLURM-USER-TOKEN} 携带 JWT；历史能力使用 slurmdb 前缀；401 自动刷新后重试一次。\\
浏览器—服务 & JSON HTTP / \code{text/event-stream} / WebSocket & 普通业务用 HTTP；聊天与安装进度用 SSE；交互终端用 WebSocket 二进制帧传 PTY 字节。\\
服务入口 & Unix Domain Socket + SSH LocalForward & Uvicorn 绑定权限受限的 Socket，浏览器经隧道访问，不在共享登录节点公开监听 TCP 端口。\\
\bottomrule
\end{longtable}

\section{关键业务流程}
四个核心流程按作业生命周期依次展开：智能对话贯穿全程，依赖流程发生在准备阶段，受控提交把准备成果落为调度实体，监控分析再把运行结果反馈回项目——最后一环回到依赖，形成闭环。

\begin{center}
\begin{tikzpicture}[node distance=5mm and 4mm, font=\small,
 b/.style={rounded corners=2mm,draw=blue,thick,fill=paper,minimum height=10mm,align=center,text width=28mm},
 a/.style={->,thick,color=muted}]
 \node[b] (prep) {依赖检查\\与安装};
 \node[b,right=of prep] (submit) {受控提交\\脚本生成与核验};
 \node[b,right=of submit] (run) {排队与运行\\状态观察};
 \node[b,right=of run] (report) {报告生成\\日志分析};
 \draw[a] (prep)--(submit); \draw[a] (submit)--(run); \draw[a] (run)--(report);
 \draw[a,bend left=30] (report.south) to node[below,font=\scriptsize]{缺失依赖写回需求} (prep.south);
\end{tikzpicture}
\end{center}

\subsection{智能对话与工具调用}
用户消息进入项目对应的 Agent 实例。模型收到 system prompt、历史消息和工具 Schema 后，可返回零个或多个工具调用；执行器将工具结果作为 \code{role=tool} 消息加入上下文，再次请求模型，直到得到自然语言回复。项目会话、文档问答和文件管理问答相互隔离，降低跨场景上下文污染。

\begin{center}
\begin{tikzpicture}[node distance=6mm, font=\small,
 n/.style={draw=blue,rounded corners,fill=paper,minimum width=39mm,minimum height=9mm,align=center},
 d/.style={diamond,draw=orange,fill=orange!10,aspect=2.4,align=center},
 a/.style={->,thick,color=muted}]
 \node[n] (msg) {接收用户消息};
 \node[n,below=of msg] (llm) {调用 Chat Completions};
 \node[d,below=of llm] (dec) {存在 tool\_calls?};
 \node[n,right=18mm of dec] (exec) {参数解析与工具执行};
 \node[n,below=of dec] (reply) {保存并流式返回回复};
 \draw[a] (msg)--(llm); \draw[a] (llm)--(dec);
 \draw[a] (dec)--node[above,font=\scriptsize]{是}(exec);
 \draw[a] (exec) |- node[pos=.25,right,font=\scriptsize]{结果回填上下文} (llm);
 \draw[a] (dec)--node[right,font=\scriptsize]{否}(reply);
\end{tikzpicture}
\end{center}

\subsection{依赖识别、确认与安装}
依赖流程不直接让模型生成并执行命令，而是遵循“确定性扫描打底、真实软件源预检、模型仅复核补充、用户最终确认”的分层设计。强调确定性优先，原因很直接：清单上的每一项最终都会变成一次真实的软件安装，一个被模型幻觉出来的包名或版本号，代价往往是几十分钟的失败排障。

\subsubsection{多源确定性扫描}
扫描器遍历项目目录，识别六类可信来源：\code{requirements.txt}、\code{environment.yml}、\code{pyproject.toml}、Pipfile、\code{setup.py} 的显式声明，以及 Shell 脚本（\code{.sh}/\code{.bash}/\code{.sbatch}）和 Markdown 文档中的安装命令与依赖清单。对 README 这类自由文本，扫描器同时识别三类内容：行内或代码块中的 \code{pip install} / \code{conda install} 命令、“Requirements/依赖”标题下的列表项，以及代码块中连续出现的“单包名+版本约束”清单（连续三行以上才采纳，避免把 Shell 片段误判为清单）；并用包名黑名单过滤 “Install xxx with conda” 这类说明性句子，防止把动词当成包名。

显式声明之外，扫描器用 AST 解析全部 Python 源码的 import 语句：先剔除标准库与项目本地模块（含 src-layout 与命名空间包），再经一张刻意保持精简的稳定别名表（如 \code{cv2} → \code{opencv-python}、\code{yaml} → \code{PyYAML}）映射为发行包名。无法归入显式依赖的 import 被标记为“待模型复核”的候选，并保留文件、行号和语句原文作为证据。用户补充的文字需求（如“需要 gromacs 2026.3”）也被确定性解析为依赖项——含用户输入的安装命令、常见软件别名表和邻近版本号提取——不依赖模型推断。

\subsubsection{版本感知预检}
预检依次回答三个问题。其一，环境里是否已装过：一次 \code{conda list} 快照，已装项默认不勾选，避免重复安装。其二，软件源里是否真的有这个包：pip 走 \code{pip index versions}，conda 走 \code{conda search --json} 并锁定渠道，结果带 10 分钟 TTL 缓存，report 与 install 两个阶段复用，避免重复下载 repodata。其三，请求的版本是否存在：按 PEP 440 规则解析版本约束（含 conda \code{version=build} 三段式的转换），与真实可用版本逐一比对。

版本不存在时，系统标记 \code{version\_mismatch} 并给出建议版本：优先取同 major.minor 的最新版，其次同 major，最后才是软件源最新版；conda 包还会展示匹配范围内最新的构建变体（如 nompi/cuda）。这一步专门防范一类高频错误：把其它集群 module 系统里的版本号（如 \code{gromacs/2019.4-gcc-9.2.0} 中的 2019.4）直接当成 conda 可安装版本。此外，conda 25.x 首次访问软件源需接受 ToS，系统在分析与安装前自动代为接受（幂等操作），用户无需感知。

\subsubsection{AI 复核的权限边界}
只有预检后仍未解决的 import 候选会交给模型分类，动作限定为外部依赖、已被显式依赖覆盖、可选、忽略、不确定五种。模型无权修改显式声明：即使模型声称某 import “已被覆盖”，系统也要求其指出的提供包确实存在于显式依赖列表中，否则保留候选——模型不能凭空删除依赖。所有 AI 建议默认不勾选，装什么、不装什么由用户在确认弹窗中最终决定。

\subsubsection{状态版本化与并发互斥}
依赖状态文件含 revision 与 analyzed\_revision。文件上传或编辑会把项目标记为 dirty，避免复用过期分析结果；缓存命中时仅做一次轻量 \code{conda list} 刷新已装状态，不重新扫描。同一项目使用线程锁与 \code{flock} 双层互斥，防止多个请求同时修改 Conda 前缀；中断产生的半成品环境会在重建前清理。

\subsection{受控作业提交}
\begin{enumerate}
  \item 用户或模型提供项目、运行子目录、命令正文和结构化资源字段。
  \item 服务校验命令不得含 shebang 或 \code{\#SBATCH}，并限制长度；校验作业名与资源参数的整数范围，并检查 CPU 分区上是否误申请了 GPU。
  \item 实时读取用户 Account/QoS 授权与分区允许范围，结合 REST 返回或静态兜底的 QoS 上限校验 CPU、GPU、内存和时限。
  \item 使用 \code{shlex} 解析逻辑命令，验证 Python/Shell 入口、配置、权重和断点路径；识别重复运行目录前缀。
  \item 后端锁定生成 \code{set -euo pipefail}、\code{cd}、日志目录和项目 Conda 激活，再调用 \code{/job/submit}。
  \item 取得 job\_id 后立即查询作业详情，从 TRES/GRES 证据核验 GPU 数；证据不足时明确标记 pending，绝不把“已提交”误报为“GPU 已确认”。
\end{enumerate}

整条提交链路如下图，橙色环节为校验网关，任一不通过即拒绝并指出具体字段：

\begin{center}
\begin{tikzpicture}[node distance=3mm and 3.5mm,>=Latex,font=\scriptsize,
 pipe/.style={rounded corners=1.5mm,draw=blue,thick,fill=paper,minimum height=10mm,align=center,text width=21mm},
 gate/.style={rounded corners=1.5mm,draw=orange,thick,fill=orange!10,minimum height=10mm,align=center,text width=21mm},
 a/.style={->,thick,color=muted}]
 \node[pipe] (in) {结构化输入\\命令正文+资源字段};
 \node[gate,right=of in] (syn) {语法与路径\\静态校验};
 \node[gate,right=of syn] (quota) {授权与配额\\Account/QoS};
 \node[pipe,right=of quota] (gen) {锁定前导生成\\cd/日志/Conda};
 \node[pipe,right=of gen] (sub) {提交\\\code{/job/submit}};
 \node[pipe,right=of sub] (verify) {TRES/GRES\\回读核验};
 \draw[a] (in)--(syn); \draw[a] (syn)--(quota); \draw[a] (quota)--(gen); \draw[a] (gen)--(sub); \draw[a] (sub)--(verify);
\end{tikzpicture}
\end{center}

\subsection{监控与结果分析闭环}
提交成功的作业写入 watch 记录。状态变为终态后，后台线程读取 Slurm 信息和 stdout/stderr 尾部。成功作业生成结果概况、关键指标与警告；失败作业生成关键错误、原因和修复步骤。若 LLM 不可用，仍生成包含原始上下文的兜底报告。系统还通过确定性正则提取 \code{ModuleNotFoundError}、\code{ImportError} 和 \code{command not found} 中的缺失依赖，写回项目需求记录，使后续依赖检查能够吸收运行反馈。

\section{功能模块设计}
下表按用户可见的模块分组，并标注各自的关键设计；其中环境与依赖、作业提交、监控与报告三个模块的实现细节，分别对应第 4 节的关键业务流程与第 9 节的技术难点。

\begin{longtable}{>{\raggedright\arraybackslash}p{29mm}>{\raggedright\arraybackslash}p{59mm}>{\raggedright\arraybackslash}p{61mm}}
\toprule
\textbf{模块} & \textbf{核心功能} & \textbf{关键设计}\\
\midrule
\endfirsthead
\toprule
\textbf{模块} & \textbf{核心功能} & \textbf{关键设计}\\
\midrule
\endhead
资源与作业看板 & 节点、CPU/GPU 状态、实时作业、个人历史作业、提交选项 & 实时数据与累计统计分口径；后端统一摘要多版本 API 字段。\\
项目工作区 & 创建项目、运行子目录、对话历史、上传、文件树 & 根路径约束、安全相对路径、内部目录保护、2 GB/1000 文件默认上传上限。\\
智能助手 & 集群查询、知识问答、文件阅读、脚本生成、日志诊断 & 19 个受限工具；实际运行数据必须工具获取；项目上下文注入。\\
模型管理 & 获取可用模型、刷新、切换 & 仅保留 deepseek/glm 家族；本地只存模型元数据和 Key 指纹，不存明文 Key。\\
环境与依赖 & 项目级 Conda、依赖扫描、选择安装、进度展示、结果验证 & Conda/Pip 命令规范化；安装互斥；SSE 分阶段反馈；分析缓存版本化。\\
模板与作业规划 & 单卡 PyTorch、多卡 DDP、CPU 批处理、Job Array、Jupyter、通用脚本 & JSON 模板参数默认值、范围校验和变量替换；支持个人模板。\\
作业提交 & Web/Agent 共用提交入口，生成脚本并调用 Slurm & 结构化资源、授权校验、路径静态检查、锁定前导、提交后资源核验。\\
监控与报告 & 终态检测、日志聚合、成功/失败报告、输出浏览和下载 & 后台分析不阻塞请求；大文件尾部预览；外部结果链接保留。\\
文件管理器 & 树状浏览、编辑、预览、下载、上传、新建、移动、复制、删除 & 路径归一化、可编辑类型与大小限制、关键内部路径保护。\\
Web 终端 & 浏览器内登录 shell、ANSI 显示、窗口缩放 & PTY 字节桥接、同源校验、最多 4 会话、SIGHUP 到 SIGKILL 的清理策略。\\
平台文档 & 文档树、正文阅读、独立问答 & 本地 Markdown 为事实来源，Embedding 检索失败时自动转关键词检索。\\
\bottomrule
\end{longtable}

\section{模型与 API 调用设计}

\subsection{使用的模型}
\begin{itemize}
  \item \textbf{默认对话模型：}\code{deepseek-v4-flash}，运行时可从当前 API Key 返回的 deepseek/glm 模型中切换。
  \item \textbf{向量模型：}\code{qwen3-embedding}，用于平台 Markdown 文档的向量化与查询向量生成。
  \item \textbf{预留重排序模型：}\code{qwen3-reranker}，配置已保留但当前检索流程未启用，文档量扩大后可接入。
\end{itemize}

\subsection{调用方式}
对话模型使用 OpenAI 兼容 Chat Completions API。请求包含 \code{model}、\code{messages}、\code{temperature}、\code{max\_tokens}，需要工具时附加 OpenAI function tool Schema。模型返回 \code{tool\_calls} 后，由 AgentLoop 执行并以对应 \code{tool\_call\_id} 回填结果。Embedding API 按 16 个片段一批请求，将向量缓存为 NumPy 文件，并以余弦相似度取最多 3 条结果，阈值为 0.3。

\begin{lstlisting}[language=Python]
client = OpenAI(api_key=os.environ["LLM_API_KEY"],
                base_url="https://api.llm.ustc.edu.cn/v1")
response = client.chat.completions.create(
    model=selected_model,
    messages=messages,
    tools=tool_definitions,
    temperature=0.3,
    max_tokens=2048,
)

vectors = client.embeddings.create(
    model="qwen3-embedding", input=batch
)
\end{lstlisting}

\subsection{工具能力矩阵}
\noindent\begin{tabularx}{\textwidth}{>{\bfseries\color{navy}}p{29mm}X p{31mm}}
\toprule
类别 & 工具能力 & 数据/执行通道\\
\midrule
作业与集群 & 实时/单作业查询、提交、取消、诊断、节点、QoS、历史 & Slurm REST\\
调度诊断 & 优先级、fairshare、实时占用、用量报表、本人授权 & 只读 Slurm CLI\\
知识与脚本 & 文档检索、模板列表、脚本生成 & 本地知识库/模板\\
项目上下文 & 文件树、按行读取文件 & 受限项目文件系统\\
\bottomrule
\end{tabularx}

\section{数据与状态设计}

系统没有引入中心数据库，而是利用用户项目目录保存与计算任务天然共址的轻量状态。项目对话、依赖分析、作业观察记录、生成脚本和报告均以 JSON、Markdown 或文本文件持久化，与代码和数据放在同一目录树下。这降低了登录节点部署复杂度，也使用户可直接检查和迁移结果。主要状态如下表：

\begin{longtable}{>{\bfseries\color{navy}}p{22mm}>{\raggedright\arraybackslash}p{76mm}>{\raggedright\arraybackslash}p{51mm}}
\toprule
\textbf{状态类别} & \textbf{存放内容} & \textbf{载体}\\
\midrule
\endfirsthead
\toprule
\textbf{状态类别} & \textbf{存放内容} & \textbf{载体}\\
\midrule
\endhead
模型配置 & base URL、Key 指纹、可用模型家族、选中模型、更新时间；不保存密钥 & JSON\\
依赖状态 & 版本号、项目 revision、最近分析 revision、结构化依赖项和报告 & 项目目录内 JSON\\
会话状态 & 按项目与运行子目录保存角色和文本消息，服务重启后可恢复 & JSON\\
作业观察 & job\_id、项目、运行目录、日志目录、最终状态、分析状态和报告位置 & JSON\\
知识索引 & 文档分块元数据与向量缓存 & JSON 与 \code{.npy}\\
\bottomrule
\end{longtable}

\section{安全性与可靠性设计}

\subsection{安全边界}
系统把安全约束拆成六层边界，模型与前端均无权跨越：

\begin{longtable}{>{\bfseries\color{navy}}p{20mm}>{\raggedright\arraybackslash}p{129mm}}
\toprule
\textbf{边界} & \textbf{约束与手段}\\
\midrule
\endfirsthead
\toprule
\textbf{边界} & \textbf{约束与手段}\\
\midrule
\endhead
密钥 & LLM Key 和 Slurm JWT 只从环境变量或本地 \code{.env} 读取，不写入代码；界面仅显示脱敏 Token 预览。\\
网络 & Unix Socket 权限配合 SSH 隧道，将 Web 访问限定在当前登录用户，不直接开放共享节点端口。\\
命令 & CLI 仅允许 \code{sprio}、\code{sshare}、\code{sstat}、\code{sreport} 和 \code{sacctmgr show/list}；argv + \code{shell=False}，参数做字符白名单与伪选项拦截。\\
文件 & 拒绝绝对路径、\code{..}、内部状态目录和越界目标；项目名清洗后再 resolve 并验证根路径前缀。\\
副作用 & 安装、提交、取消由用户明确触发；Agent 提交必须绑定项目级 handler，不能从命令行默认会话绕过。\\
资源 & 结构化整数范围、账户/分区/QoS 交叉授权和限额校验；GPU 进入 REST job description 而非只写脚本注释。\\
\bottomrule
\end{longtable}

\subsection{可靠性机制}
Slurm 请求设置统一超时并规范化 401/404/网络异常；401 最多刷新 JWT 后重试一次。知识检索、结果分析和 QoS 限额读取均设计了降级路径。长耗时的 Conda 初始化、依赖安装和作业报告采用后台执行或 SSE，不阻塞主请求。终端断连时先发送 SIGHUP，宽限后强制终止，避免僵尸 shell。多处输出设置大小或条数上限，控制浏览器、服务和模型上下文的资源消耗。

\section{技术难点与解决方案}

\subsection{异构 Slurm 接口语义对齐}
Slurm REST、slurmdbd 和 CLI 的字段、可用性与统计口径并不完全一致。例如 v0.0.41 的 partition 查询参数可能被忽略，系统在客户端二次过滤；提交端点使用单数 \code{/job/submit}，nodes 需字符串类型，time limit 字段和 environment 也有平台特定要求。系统把这些差异封装在 SlurmClient 中，并以 CLI 只读能力补足 REST 空白，避免上层业务同时理解多套语义。

\subsection{模型生成与真实执行之间的可信鸿沟}
模型能够生成看似正确但路径不存在、资源字段缺失或夹带 \code{\#SBATCH} 的命令。项目采用“双阶段生成”：模型只给正文，后端锁定关键部分；再进行命令语法和文件存在性检查。提交成功也不等于资源申请生效，因此系统查询作业详情并解析 TRES/GRES 作为证据。这使“自然语言生成”具备可验证的落点。

\subsection{依赖安装引擎的失败分类与分层自愈}
依赖检查解决了“装什么”，安装引擎要解决“装失败之后怎么办”。科学计算依赖的安装失败是常态而非例外：版本号在软件源中不存在、求解器冲突、源码包缺编译器、大体积下载中断，每一类失败的正确处置方式都不同。系统没有把这些判断全部交给大模型，而是先让确定性分类器给出结论，模型只承担剩余的少数疑难情形。

\textbf{PTY 执行与停滞判定。}每条安装命令在 PTY 中执行，与 Web 终端共用同一套机制：conda/pip 自带的下载进度条、颜色和交互提示原样经 SSE 推送到浏览器中的 xterm.js，用户看到的就是真实执行过程，不存在模拟进度。超时判定采用“停滞”语义而非固定总时长——PyTorch 这类包要下载数 GB，慢速网络下健康下载也会持续一小时以上；连续十分钟无任何新输出才判定卡死（进度条刷新即是存活的证据），另设四小时总时长兜底，两者均可通过环境变量调整。超时不进入失败分类器、不喂给模型：模型修不了“时间不够”，直接终止并明示原因，重试会利用 pip/conda 缓存显著加快。

\textbf{失败分类器。}命令失败时，分类器在输出尾部 60 行内做锚点匹配，把失败归为八类：磁盘/权限、网络、包不存在、构建失败、求解冲突、版本不存在、Python 版本不匹配、未归类，并提取缺失包名、冲突上下文、HTTP URL 等结构化信息。以尾部为主要匹配面，是为了避开 pip 构建日志中部的网络重试警告，防止把健康的重试过程误判为网络故障。

\textbf{分层处置策略。}处置按“确定性优先”排序：
\begin{enumerate}
  \item \textbf{确定性自动修复：}求解冲突和版本不存在先自动去掉版本约束重试一次，让 conda/pip 求解器自选兼容组合，全程不打断用户。原失败命令被标记为“已被取代”，最终成败只看重试结果，不因历史失败误报。
  \item \textbf{直接终止：}磁盘、网络、包不存在、构建失败和超时属于模型修不了的问题（模型无法凭空变出磁盘空间或编译器），系统给出精确原因立即结束，不浪费一轮模型调用，也不让模型猜测不存在的包名。
  \item \textbf{LLM 兜底修正，人工确认后执行：}其余失败（版本残留、Python 版本不匹配、未归类）交给模型。系统把真实软件源查询结果（可用版本与构建列表）、集群硬件上下文和目标环境 Python 版本一并注入提示词，要求模型只输出修正后的安装命令，版本号必须来自查询结果；服务端再做一道确定性校验——若模型给 py3.10 环境选了 py311 构建，整体拒绝。修正方案以弹窗形式提交用户确认：确认前安装流程挂起等待，终止安装操作也能唤醒该等待，消除“用户在弹窗出现前点了终止”的竞态。
\end{enumerate}

三层策略的入口与边界如下图所示——分类结果决定走哪条路，模型只出现在第三条：

\begin{center}
\begin{tikzpicture}[font=\scriptsize,
 head/.style={rounded corners=1.5mm,draw=orange,thick,fill=orange!10,minimum height=9mm,align=center,text width=40mm},
 cond/.style={rounded corners=1.5mm,draw=muted,thick,fill=paper,minimum height=10mm,align=center,text width=40mm},
 act/.style={rounded corners=1.5mm,draw=blue,thick,fill=blue!6,minimum height=10mm,align=center,text width=56mm},
 stop/.style={rounded corners=1.5mm,draw=red!60!black,thick,fill=red!8,minimum height=10mm,align=center,text width=56mm},
 a/.style={->,thick,color=muted}]
 \node[head] (cls) at (0,0) {失败分类器\\输出尾部 60 行锚点匹配};
 \node[cond] (c1) at (-3.3,-1.8) {求解冲突 / 版本不存在};
 \node[act]  (a1) at (4.6,-1.8) {自动去版本约束重试一次，不打断用户};
 \node[cond] (c2) at (-3.3,-3.2) {磁盘 / 网络 / 包不存在 / 构建失败 / 超时};
 \node[stop] (a2) at (4.6,-3.2) {直接终止，给出精确原因，不调用模型};
 \node[cond] (c3) at (-3.3,-4.6) {版本残留 / Python 版本不匹配 / 未归类};
 \node[act]  (a3) at (4.6,-4.6) {LLM 兜底修正 → 弹窗人工确认后执行};
 \coordinate (hub) at ($(cls.south)+(0,-0.35)$);
 \draw[a] (cls.south) -- (hub) -| (c1.north);
 \draw[a] (hub) -| (c2.north);
 \draw[a] (hub) -| (c3.north);
 \draw[a] (c1)--(a1); \draw[a] (c2)--(a2); \draw[a] (c3)--(a3);
\end{tikzpicture}
\end{center}

\textbf{终止与残留防护。}无论超时还是用户手动终止，进程都不直接 SIGKILL：先发送 SIGINT/SIGHUP 让 pip 停止写入 site-packages、conda 中止事务并标记未解压包，等待六秒后仍存活才强制终止，降低留下残缺环境的概率。中断的安装支持增量重试：已成功执行的命令可跳过，只重跑失败及未执行部分。

\textbf{安装后验证与反馈闭环。}全部命令结束后，系统用一次 \code{conda list} 快照逐一核验勾选项是否真的进入目标环境，而非以命令退出码为准。作业失败日志中被确定性抽取的缺失依赖会写回项目需求记录，参与下一次“检查依赖”分析，形成安装层面的运行反馈闭环。

\subsection{共享登录节点上的环境与并发}
服务环境、项目计算环境、Conda base 必须隔离；同时初始化同一项目可能破坏前缀。系统为每个项目创建独立前缀，生成自适应 \code{activate.sh}，用线程锁覆盖 FastAPI 线程池并用 \code{flock} 覆盖多进程，且对中断半成品进行自愈处理。

\subsection{跨生命周期上下文闭环}
一个作业可能排队数小时，HTTP 请求早已结束。系统将作业注册为可观察记录，独立追踪终态，再把报告写回原项目/运行目录会话。失败日志中确定性抽取的缺失依赖还会反馈到下一次依赖分析，形成“准备 - 提交 - 运行 - 诊断 - 修正”的闭环。

\subsection{浏览器终端的系统编程细节}
普通 subprocess 管道不能正确承载交互式 shell、ANSI 控制和窗口大小变化。系统通过 PTY 提供真实终端语义，在 asyncio 与阻塞读取线程之间安全转发字节，并处理 SIGWINCH、Origin、并发配额和断连后的进程回收。这一模块涉及浏览器协议、异步 I/O 和 Unix 进程控制的协同。

\section{创新点}

\begin{enumerate}
  \item \textbf{模型与执行权解耦的受控提交。}模型只产出意图和命令正文，后端锁定所有关键调度字段、环境与日志，并在提交后用 TRES/GRES 证据反向核验资源申请是否生效。
  \item \textbf{面向具体平台的模型校准。}系统提示词内置在该集群实测确认的平台事实（分区名称区分大小写、分区—账户—QoS 对应关系、默认配额、scrontab 禁用）与数据口径纪律（当前队列以 \code{list\_jobs} 实时逐条统计为准，\code{get\_diag} 的累计统计禁止混作当前状态），使模型回答贴合本平台实际而非通用 HPC 印象。
  \item \textbf{带证据链的知识检索。}平台文档按标题分块并向量化（磁盘缓存、低于相似度阈值丢弃），回答保留来源与相似度；向量服务不可用时降级为关键词检索，构成“证据优先”原则在知识侧的落地。
  \item \textbf{版本感知的依赖预检与失败分类自愈。}依赖在安装前先与真实软件源逐项比对，针对“把其它集群 module 版本号当 conda 版本”的高频错误按同 major.minor → 同 major → 最新版的三级回退给出建议版本；安装失败经确定性分类器归为八类，常见冲突自动去版本重试，环境类失败直接终止并明示原因，仅少数疑难情形交由模型修正且必须经用户确认；超时判定采用“停滞”语义而非固定总时长。
  \item \textbf{运行反馈驱动的环境闭环。}作业失败日志中确定性抽取的缺失依赖写回项目需求；多运行目录共享项目依赖而会话隔离，报告自动回流原会话，使环境知识随实际使用逐步补全。
\end{enumerate}

\section{部署与运行设计}

\subsection{运行环境}
要求 Python 3.10+、Conda、可访问 Slurm REST API 与 OpenAI 兼容模型服务，并部署在可执行 Slurm CLI 的登录节点。服务端依赖主要包括 FastAPI、Uvicorn、OpenAI SDK、Requests、NumPy、Pydantic、WebSockets、ptyprocess 与 Packaging。前端为单页原生 Web，不需要 Node 构建链。

\subsection{启动与访问链路}
\begin{lstlisting}[language=bash]
# remote login node
python3 -m venv .venv
source .venv/bin/activate
pip install -r requirements.txt
export LLM_API_KEY="..."
./start-server.sh

# local ~/.ssh/config
Host 107.ustc.edu.cn
  LocalForward 8080 /home/<user>/slurm-ai-agent/server.sock
\end{lstlisting}
本地浏览器访问 \code{http://localhost:8080}，请求经 SSH 隧道转发到远端 Unix Socket。健康检查可通过 \code{curl --unix-socket server.sock http://localhost/health} 完成。

\section{测试与验收建议}

\begin{longtable}{>{\raggedright\arraybackslash}p{31mm}>{\raggedright\arraybackslash}p{68mm}>{\raggedright\arraybackslash}p{50mm}}
\toprule
\textbf{测试域} & \textbf{关键用例} & \textbf{通过标准}\\
\midrule
功能 & 创建项目、上传、依赖检查/安装、生成正文、提交、监控、查看报告 & 主流程无断点且状态可恢复\\
资源校验 & 无授权账户、QoS 越界、CPU 分区申请 GPU、超时限 & 提交前明确拒绝并指出字段\\
路径安全 & 绝对路径、\code{..}、重复运行目录、缺失入口/配置 & 越界拒绝，常见路径错误可定位到行\\
模型工具 & 多工具连续调用、工具异常、无工具纯问答 & 上下文结构合法，错误不被掩盖\\
Slurm 兼容 & JWT 过期、服务端忽略过滤参数、GPU TRES 延迟可见 & 单次刷新重试，客户端过滤，核验状态准确\\
依赖并发 & 同项目并行初始化/安装、中断半成品 & 串行化成功且可自愈\\
终端 & Origin 不匹配、超过 4 会话、调整尺寸、异常断连 & 拒绝策略与进程清理符合预期\\
降级 & Embedding 不可用、LLM 报告失败 & 关键词结果或原始日志报告仍可用\\
\bottomrule
\end{longtable}

\section{可演进方向}
系统仍有四条明确的演进主线。

\textbf{智能体交互的流式化。}当前智能体在一个工具调用循环结束后才整体回显结果，长流程下的等待感明显。后续可借鉴 DeepSeek 等前沿系统在 harness 层的流式设计，把回复文本、工具调用意图与工具执行进度逐步流式回显到前端：模型边生成边呈现，工具执行状态实时可见，既缩短感知延迟，也让“模型正在做什么”对用户全程透明。这一改动仅涉及展示层与事件通道，不改变既有的工具权限与校验边界。

\textbf{依赖选择的可视化。}当前依赖确认是一份平铺的勾选清单，能表达“装不装”，但难以表达“有哪些可选”。预检阶段其实已经拿到了真实软件源的结构化数据（渠道、版本列表、构建变体），后续可把它转化为可视化的选择界面：以包为根节点，向外发散展示不同 channel（conda-forge、bioconda 等）、不同版本与构建（nompi/cuda、py310/py311）的可选分支，让“同名包在不同渠道的差异”“版本与构建的兼容关系”一目了然；交互形式既可以保留命令行式的补全与筛选，也可以做成发散图浏览。由于数据全部来自现有 \code{conda search}/\code{pip index} 预检结果，该改动只增强呈现与选择方式，不扩大执行边界，最终装什么仍由用户逐项确认。

\textbf{终端便捷性的持续提升。}Web 终端目前提供的是“可用”的登录 shell，后续可向“好用”演进。一是常在与可达性：会话与 PTY 解耦持久化（tmux 式），浏览器关闭或网络中断后重新打开即可恢复现场，长时运行的交互进程不因页面刷新而丢失。二是输入便捷化：借鉴 fish 等现代 shell 的体验，基于命令历史提供灰字建议、参数补全与语法高亮。三是与超算场景特征绑定：补全候选不只来自通用历史，还结合系统已掌握的上下文——Slurm 命令的常用参数、本项目路径与 Conda 环境名、个人历史作业的提交命令与模板变量，让“上次是怎么提交的”一次按键即可找回。这些增强均发生在用户自己的 shell 会话内，不改变终端的权限边界与并发限制。

\textbf{面向更多学科场景的验证扩展。}当前端到端流程已在深度学习训练等场景验证，后续将系统性地跑通更多应用场景：覆盖物理、化学、材料、生物等学科的常用计算软件与公开数据集（如 GROMACS 分子动力学、第一性原理计算、Job Array 批量任务等），为每类场景沉淀模板、依赖清单与验收用例，形成可回归的端到端测试集。这既能暴露各学科软件在依赖求解、许可文件与并行构建上的差异，也为依赖别名库、模板库和报告指标的完善提供真实输入。

在此基础上，其余工程改进包括：自动化单元/集成测试与 OpenAPI 契约测试；为作业监控引入更稳定的任务队列；在文档规模扩大后接入 reranker；建立依赖包名与 import 名映射库；对报告增加资源效率指标；对高风险文件操作增加回收站机制；在不扩大模型权限的前提下，将受控工具扩展到更多平台能力。

\section{结论}
Slurm AI Agent 的价值不只是为 Slurm 加一个聊天界面，而是围绕 HPC 作业生命周期建立了一套“语义理解可智能、执行边界可确定、运行结果可追溯”的工程体系。通过项目化状态、Function Calling、REST/CLI 分层、受控脚本生成、资源回读核验和反馈式依赖修正，系统在易用性、安全性与真实性之间取得平衡，适合教学、科研训练和批量实验等共享算力场景。

\end{document}
